% ============================================================
%  part7/ch41-philosophical-consequences.tex
%  Chapter 41: The Monoturn Theorem — Closing Synthesis
% ============================================================

\chapter{The Monoturn Theorem and Closing Synthesis}
\label{ch:philosophical-consequences}

\epigraph{%
  Reality is not what is observed.\\
  Reality is what remains admissibly reconstructible\\
  from observation.%
}{---}

\begin{WhyChapter}
  Every chapter of this book has studied the relationship
  between an observation and the structure that generated it.
  This chapter brings those threads together.
  The central diagram has appeared in many forms:
  maps, photographs, detectors, embeddings, horizons,
  sheaves, and RSVP field configurations.
  Here it appears for the final time in its most precise form.
\end{WhyChapter}

\section{The Monoturn theorem}

\begin{theorem}[Monoturn reconstruction]
  \label{thm:monoturn-final}
  \statusproven{}
  Let $\mathcal{A}$ be an admissibility sheaf over $X$.
  If (1) $\check{\delta}s = 0$, (2) $H^1(X,\mathcal{A}) = 0$,
  and (3) $E(X) < \infty$, then a unique admissible global
  reconstruction exists.
\end{theorem}

\begin{proof}
  Proved in Chapter~\ref{ch:monoturn} as
  Theorem~\ref{thm:holonomic-recon}.
\end{proof}

When $H^1 \neq 0$, the Monoturn exceeds reconstruction.
That excess is not mystical. It is a measurable obstruction class.

\section{Examples of the same structure}

\begin{WorkedExample}[Photograph]
  A photograph records shape, brightness, texture.
  It loses depth, temperature, chemical composition.
  The photograph is not false. It is a projection.
  $\ker(\pi) \neq 0$: depth, temperature, history lie in the kernel.
\end{WorkedExample}

\begin{WorkedExample}[Map]
  A map preserves adjacency and routes.
  It loses elevation, vegetation, weather.
  $\ker(\pi) \neq 0$: terrain and weather lie in the kernel.
\end{WorkedExample}

\begin{WorkedExample}[Cosmic horizon]
  An observer inside horizon $H_i$ sees only $H_i$.
  The Monoturn is $\mathcal{M} = \bigcup_i H_i$
  yet $\mathcal{M} \neq H_i$.
  The horizon is not an observational defect.
  It is a structural feature of reconstruction.
  $\ker(\pi_O \circ F_O) \neq 0$: the beyond-horizon structure
  lies in the kernel.
\end{WorkedExample}

\section{The closing diagram}

\[
  \mathcal{M}
  \xrightarrow{F}
  \mathcal{A}
  \xrightarrow{\pi}
  \mathcal{O},
  \qquad
  \ker(\pi \circ F) \neq 0.
\]

\begin{proposition}[Observable simplicity $\neq$ ontological simplicity]
  \label{prop:obs-onto}
  \statusproven{}
  If $\ker(\pi \circ F) \neq 0$, then observational equivalence
  does not imply ontological identity:
  there exist $\mathcal{M}_1 \neq \mathcal{M}_2$ with
  $(\pi \circ F)(\mathcal{M}_1) = (\pi \circ F)(\mathcal{M}_2)$.
\end{proposition}

\begin{proof}
  Non-injectivity gives pairs that are observationally
  indistinguishable but ontologically distinct.
  This was proved in general in Part~I and applies here to
  the cosmological case.
\end{proof}

\section{Mereological decomposition and reconstruction}

The forward Mereological Space Ontology (MSO) decomposes:
\[
  \mathcal{M}
  \to \text{Filaments}
  \to \text{Galaxies}
  \to \text{Planets}
  \to \text{Life}
  \to \text{Cells}
  \to \text{Atoms}
  \to \text{Fields}.
\]

The reverse MSO reconstructs:
\[
  \text{Observer}
  \to \text{Fields}
  \to \text{Matter}
  \to \text{Life}
  \to \text{Planet}
  \to \text{Galaxy}
  \to \mathcal{M}.
\]

These are not inverses. They are adjoint operations:
$F \dashv G$.
The forward direction asks what something is made of.
The reverse asks what something must be part of.

The Monoturn is the fixed point:
$\mathcal{M} \simeq G(F(\mathcal{M}))$
when $H^1(\mathcal{M}, \mathcal{A}) = 0$.

\section{The final statement}

\principle{
  Reality is not what is observed.\\[4pt]
  Reality is what remains admissibly reconstructible\\[4pt]
  from observation.
}

Mathematically:
\[
  \boxed{\mathcal{M} \not\cong \pi(F(\mathcal{M})).}
\]

The entire monograph is the study of that difference:
what lies in $\ker(\pi \circ F)$, how much can be recovered
from what remains, and what the structure of the obstruction
reveals about the whole.

\begin{HolonomyNote}
  After traversing the entire book, the holonomy is now visible.

  Part~I showed that every observation is a projection.
  Part~II showed that reconstruction is imperfect whenever
  the projection is non-injective.
  Part~III showed that physical systems evolve toward simpler
  descriptions through coarsening.
  Part~IV gave the RSVP field theory its analytic foundations.
  Part~V formalised admissibility as a sheaf condition and
  obstructions as cohomology.
  Part~VI stated what is proven, conjectured, and open.
  Part~VII applied the framework to the observable universe.

  The invariant that becomes visible only after the full
  traversal: every chapter has been studying the kernel of a
  projection.

  Maps, photographs, memories, detectors, embeddings, horizons,
  sheaves, defects, coarsening domains, cosmic voids, and
  causal horizons are all instances of the same object:
  $\ker(\pi \circ F)$.

  The Holonomic Space is the space of all such kernels,
  equipped with the structure that determines what can be
  reconstructed from what remains.
\end{HolonomyNote}

\begin{ProblemSet}
\problem{
  State the five conditions under which the Monoturn is
  reconstructible from a finite observer cover.
  For each condition, give an example of a physical situation
  where it fails.
}
\problem{
  \textbf{Inverse problem.}
  Given observations $\mathcal{O} = \pi(F(\mathcal{M}))$,
  describe the admissibility class $[\mathcal{M}]_\pi
  = \{N : \pi(F(N)) = \mathcal{O}\}$.
  Under what conditions is this class a singleton?
}
\problem{
  The forward MSO is a functor $F : \mathbf{Glob} \to \mathbf{Loc}$.
  The reverse MSO is $G : \mathbf{Loc} \to \mathbf{Glob}$.
  Show that $F \dashv G$ implies $G \circ F \neq \mathrm{id}$
  in general.
  What does $G(F(\mathcal{M})) \neq \mathcal{M}$ mean
  physically for a cosmological observer?
}
\problem{
  Compute $H^1(S^1, \mathbb{Z})$ using the Mayer--Vietoris
  sequence.
  Interpret the result in terms of the Monoturn:
  what global structure of the circle cannot be reconstructed
  from any proper open subcover?
}
\problem{
  \textbf{Open research problem.}
  Suppose $H^1(\mathcal{M}, \mathcal{A}^{(k)}_{\mathrm{RSVP}}) \neq 0$
  for the cosmological admissibility sheaf.
  Propose an observational signature that would indicate
  the presence of nontrivial obstruction classes.
  (There is no known answer. This is an open problem.)
}
\end{ProblemSet}
